% Source: https://artofproblemsolving.com/community/c3h1338687_mathcounts_2015_state_22
% Stable identity: aops-topic-1338687-q22
A large bar ($8$ kg in ratio $4:1:3$) contains $4$ kg iron, $1$ kg nickel, $3$ kg lead. A medium bar ($6$ kg in ratio $2:1:3$) contains $2$, $1$, $3$ kg. A small bar ($3$ kg in ratio $1:1:1$) contains $1$, $1$, $1$ kg. With $L$ large, $M$ medium and $s$ small bars:
\[
\begin{aligned}
\text{iron: }&4L+2M+s=40,\\
\text{nickel: }&L+M+s=20,\\
\text{lead: }&3L+3M+s=40.
\end{aligned}
\]
Subtracting the nickel equation from the lead equation gives $2L+2M=20$, so $L+M=10$ and $s=10$. Then the iron equation gives $4L+2M=30$, so $2L+M=15$, $L=5$, $M=5$. All are non-negative integers, so she uses $\boxed{10}$ small bars (with $5$ large and $5$ medium bars).
